% Part:sets-functions-relations
% Chapter: sets
% Section: reduction

\documentclass[../../../include/open-logic-section]{subfiles}

\begin{document}

\olfileid{sfr}{siz}{red}

\olsection{Reductie}

\begin{editorial}
  Deze paragraaf bewijst overaftelbaarheid met een reductie en sluit aan bij
  de resultaten in \olref[nen]{sec}. Een andere, iets kortere versie die
  aansluit bij de resultaten in
  \olref[nen-alt]{sec}, staat in \olref[red-alt]{sec}.
\end{editorial}

Met een diagonaalargument hebben we aangetoond dat $\Pow{\PosInt}$
!!{nonenumerable} is. We hadden al een bewijs dat $\Bin^\omega$, de
verzameling van alle oneindige rijen van $0$'en en $1$'en,
!!{nonenumerable} is. We kunnen ook op een andere manier bewijzen dat
$\Pow{\PosInt}$ !!{nonenumerable} is: toon aan dat \emph{als
$\Pow{\PosInt}$ !!{enumerable} is, dan $\Bin^\omega$ eveneens
!!{enumerable} is}. Omdat we weten dat $\Bin^\omega$ niet
!!{enumerable} is, kan $\Pow{\PosInt}$ dat evenmin zijn. Dit heet het
ene probleem tot het andere \emph{reduceren}---in dit geval reduceren we
het probleem om $\Bin^\omega$ op te sommen tot het probleem om
$\Pow{\PosInt}$ op te sommen. Een oplossing van het laatste---een
opsomming van $\Pow{\PosInt}$---zou een oplossing van het eerste---een
opsomming van $\Bin^\omega$---opleveren.

Hoe reduceren we het probleem om een verzameling~$B$ op te sommen tot
het probleem om een verzameling~$A$ op te sommen? We geven een manier
om een opsomming van~$A$ om te zetten in een opsomming van~$B$. Het
eenvoudigst is een surjectieve functie $f\colon A \to B$ te
definiëren. Als $x_1$, $x_2$, \dots{} een opsomming van~$A$ vormen,
dan zouden $f(x_1)$, $f(x_2)$, \dots{} de verzameling~$B$ opsommen. In
ons geval zoeken we een surjectieve functie $f\colon \Pow{\PosInt} \to
\Bin^\omega$.

\begin{prob}
Toon aan dat als er een injectieve functie $g\colon B \to A$ bestaat
en $B$~!!{nonenumerable} is, ook~$A$ !!{nonenumerable} is. Laat daartoe
zien hoe met~$g$ een opsomming van~$A$ kan worden omgezet in een
opsomming van~$B$.
\end{prob}

\begin{proof}[Bewijs van {\olref[nen]{thm:nonenum-pownat}} door reductie]
Stel dat $\Pow{\PosInt}$ !!{enumerable} was en dat er dus een
opsomming van bestond, $Z_{1}$, $Z_{2}$, $Z_{3}$, \dots

Definieer de functie $f \colon \Pow{\PosInt} \to \Bin^\omega$ door
$f(Z)$ de rij $s_{k}$ te laten zijn waarvoor $s_{k}(n) = 1$ precies
dan als $n \in Z$, en waarvoor $s_k(n) = 0$ in alle andere gevallen.
Dit definieert inderdaad een functie, want als $Z \subseteq \PosInt$,
is elke $n \in \PosInt$ ofwel !!a{element} van $Z$, ofwel niet. Zo
wordt bijvoorbeeld de verzameling $2\PosInt = \{2, 4, 6, \dots\}$ van
positieve even getallen afgebeeld op de rij $010101\dots$, de lege
verzameling op $0000\dots$ en de verzameling $\PosInt$ zelf op
$1111\dots$.

De functie is ook !!{surjective}: elke rij van $0$'en en $1$'en bepaalt
een verzameling positieve gehele getallen, namelijk de indexen waarop de
rij een~$1$ heeft. Meer precies, neem $s \in \Bin^\omega$. Definieer
$Z \subseteq \PosInt$ door:
\[
Z = \Setabs{n \in \PosInt}{s(n) = 1}
\]
Dan geldt $f(Z) = s$, zoals uit de definitie van~$f$ volgt.

Beschouw nu de lijst
\[
f(Z_1), f(Z_2), f(Z_3), \dots
\]
Omdat $f$ !!{surjective} is, moet elk element van $\Bin^\omega$ voor
enig argument als waarde van~$f$ optreden en dus in de lijst voorkomen.
Deze lijst somt daarom heel~$\Bin^\omega$ op.

Als $\Pow{\PosInt}$ dus !!{enumerable} was, zou $\Bin^\omega$
!!{enumerable} zijn. Maar $\Bin^\omega$ is !!{nonenumerable}
(\olref[nen]{thm:nonenum-bin-omega}). Dus is $\Pow{\PosInt}$
!!{nonenumerable}.
\end{proof}

\begin{explain}
De richting van een reductie kan gemakkelijk verwarrend zijn. Zo toont
een surjectieve functie $g \colon
\Bin^\omega \to B$ 
\emph{niet} aan dat $B$ !!{nonenumerable} is. (Beschouw $g \colon
\Bin^\omega \to \Bin$ gedefinieerd door $g(s) = s(1)$, de functie die
een rij van $0$'en en $1$'en afbeeldt op haar eerste !!{element}. Deze
functie is !!{surjective}, omdat sommige rijen met $0$ beginnen en
andere met $1$. Maar $\Bin$ is eindig.) Merk ook op dat de functie~$f$
!!{surjective} moet zijn, want anders gaat het argument niet op:
$f(x_1)$, $f(x_2)$, \dots{} zouden dan niet noodzakelijk alle
!!{element}en van~$B$ bevatten. Zo definieert bijvoorbeeld
\[
h(n) = \underbrace{000\dots0}_{\text{$n$ keer $0$}}
\]
een functie $h\colon \PosInt \to \Bin^\omega$, maar $\PosInt$ is
!!{enumerable}.
\end{explain}

\begin{prob}
Toon met een reductieargument aan dat de verzameling van alle
\emph{verzamelingen van} paren positieve gehele getallen
!!{nonenumerable} is.
\end{prob}

\begin{prob}\label{sfr:siz:red:prob:nat-nat}
  Toon met een reductieargument aan dat de verzameling~$X$ van alle
  functies $f\colon \Nat \to \Nat$ !!{nonenumerable} is (Aanwijzing:
  geef een surjectieve functie van $X$ naar~$\Bin^\omega$.)
\end{prob}

\begin{prob}
Toon met een reductieargument aan dat $\Nat^\omega$, de verzameling van
oneindige rijen natuurlijke getallen, !!{nonenumerable} is.
\end{prob}

\begin{prob}
Laat $P$ de verzameling zijn van functies van de verzameling positieve
gehele getallen naar de verzameling $\{0\}$, en laat $Q$ de verzameling
zijn van \emph{partiële} functies van de verzameling positieve gehele
getallen naar de verzameling $\{0\}$. Toon aan dat $P$~!!{enumerable}
is en $Q$~niet. (Aanwijzing: reduceer het probleem om $\Bin^\omega$ op
te sommen tot het probleem om~$Q$ op te sommen.)
\end{prob}

\begin{prob}
Laat $S$ de verzameling zijn van alle surjectieve functies van de
verzameling positieve gehele getallen naar de verzameling \{0,1\}, dat
wil zeggen: $S$ bestaat uit alle surjectieve functies~$f\colon
\PosInt \to \Bin$. Toon aan dat $S$ !!{nonenumerable} is.
\end{prob}

\begin{prob}
Toon aan dat de verzameling~$\Real$ van alle reële getallen
!!{nonenumerable} is.
\end{prob}
\end{document}
